% GENERATED FILE. Edit canonical lessons, then run npm run generate:books.
% canonical-source-hash: fnv1a64:3335abf8ae1e54fd
% canonical-lessons: JA-C146-nemuru, JA-C146-gussuri, JA-C146-sameru, JA-C146-okosu, JA-C146-suimin

\chapter{Sleeping Soundly and Waking Up}
\label{ch:ja-sleeping-soundly-and-waking-up}
\hlchaptermodality{146}

\begin{chapteropening}
\textbf{By the end of this chapter:} \emph{I can talk about how I sleep and wake up.}

Complete the last lesson of chapter 146: I can talk about how I sleep and wake up.
\end{chapteropening}

\section[\texorpdfstring{\ja{ねむる} (nemuru)}{ねむる (nemuru)}]{\ja{ねむる} (nemuru) — to fall asleep; to sleep}

\label{lesson:JA-C146-nemuru}



\begin{quote}
Before the new one: say the Japanese for open for business, then the Japanese for not yet open.
\end{quote}

\subsection*{You'll want to know: \ja{ねむる}}
\textbf{\ja{ねむる}} — \emph{nemuru} — \textquotedblleft{}to fall asleep; to sleep\textquotedblright{}.

\begin{quote}
\textbf{\ja{こどもは} \ja{よく} \ja{ねむります。}} — \emph{kodomo wa yoku nemurimasu} — \textquotedblleft{}The child sleeps well.\textquotedblright{}
\end{quote}

Polite \textbf{\ja{ねむります}}, \emph{nemurimasu}; a godan verb.

\begin{cousinweb}
\textbf{\ja{ねむる}} (\emph{nemuru}) is a native Japanese word (wago), inherited rather than borrowed from Chinese; its kanji was attached later as a kun'yomi, and its older origin is not settled (JMdict, KANJIDIC2).
\end{cousinweb}

\begin{grammarlens}[title={What you've built}]
One more word for everyday life at A2.
\end{grammarlens}

\subsection*{Your turn}
\begin{itemize}
\raggedright
  \item \emph{Say it:} \emph{nemuru}
  \item \emph{Say it:} \emph{nemuru}, once more
  \item \emph{Say it:} \emph{kodomo wa yoku nemurimasu} — The child sleeps well.
  \item \emph{Recall:} say the Japanese for open for business, then the Japanese for not yet open, then say \emph{nemuru} again
\end{itemize}

\begin{tcolorbox}[breakable,colback=teal!4,colframe=teal!35!black,title={Before you move on}]
What does \ja{ねむる} mean? (\textquotedblleft{}To fall asleep; to sleep\textquotedblright{}.) Say it once more.
\end{tcolorbox}
\section[\texorpdfstring{\ja{ぐっすり} (gussuri)}{ぐっすり (gussuri)}]{\ja{ぐっすり} (gussuri) — soundly (asleep)}

\label{lesson:JA-C146-gussuri}



\begin{quote}
Before the new one: say the Japanese for not yet open, then the Japanese for to fall asleep.
\end{quote}

\subsection*{You'll want to know: \ja{ぐっすり}}
\textbf{\ja{ぐっすり}} — \emph{gussuri} — \textquotedblleft{}soundly (asleep)\textquotedblright{}.

\begin{quote}
\textbf{\ja{まいにち} \ja{ぐっすり} \ja{ねむります。}} — \emph{mainichi gussuri nemurimasu} — \textquotedblleft{}I sleep soundly every day.\textquotedblright{}
\end{quote}

\begin{cousinweb}
\textbf{\ja{ぐっすり}} (\emph{gussuri}) is a native Japanese mimetic word, inherited rather than borrowed from Chinese, and its older origin is not settled (JMdict).
\end{cousinweb}

\begin{grammarlens}[title={What you've built}]
One more word for everyday life at A2.
\end{grammarlens}

\subsection*{Your turn}
\begin{itemize}
\raggedright
  \item \emph{Say it:} \emph{gussuri}
  \item \emph{Say it:} \emph{gussuri}, once more
  \item \emph{Say it:} \emph{mainichi gussuri nemurimasu} — I sleep soundly every day.
  \item \emph{Recall:} say the Japanese for not yet open, then the Japanese for to fall asleep, then say \emph{gussuri} again
\end{itemize}

\begin{tcolorbox}[breakable,colback=teal!4,colframe=teal!35!black,title={Before you move on}]
What does \ja{ぐっすり} mean? (\textquotedblleft{}Soundly (asleep)\textquotedblright{}.) Say it once more.
\end{tcolorbox}
\section[\texorpdfstring{\ja{さめる} (sameru)}{さめる (sameru)}]{\ja{さめる} (sameru) — to wake up; to come round}

\label{lesson:JA-C146-sameru}



\begin{quote}
Before the new one: say the Japanese for to fall asleep, then the Japanese for soundly.
\end{quote}

\subsection*{You'll want to know: \ja{さめる}}
\textbf{\ja{さめる}} — \emph{sameru} — \textquotedblleft{}to wake up; to come round\textquotedblright{}.

\begin{quote}
\textbf{\ja{あさ} \ja{めが} \ja{さめます。}} — \emph{asa me ga samemasu} — \textquotedblleft{}I wake up in the morning.\textquotedblright{}
\end{quote}

Polite \textbf{\ja{さめます}}, \emph{samemasu}; an ichidan verb. With \textbf{\ja{め}}, \emph{me}, it means one's eyes open from sleep.

\begin{cousinweb}
\textbf{\ja{さめる}} (\emph{sameru}) is a native Japanese word (wago), inherited rather than borrowed from Chinese; its kanji was attached later as a kun'yomi, and its older origin is not settled (JMdict, KANJIDIC2).
\end{cousinweb}

\begin{grammarlens}[title={What you've built}]
One more word for everyday life at A2.
\end{grammarlens}

\subsection*{Your turn}
\begin{itemize}
\raggedright
  \item \emph{Say it:} \emph{sameru}
  \item \emph{Say it:} \emph{sameru}, once more
  \item \emph{Say it:} \emph{asa me ga samemasu} — I wake up in the morning.
  \item \emph{Recall:} say the Japanese for to fall asleep, then the Japanese for soundly, then say \emph{sameru} again
\end{itemize}

\begin{tcolorbox}[breakable,colback=teal!4,colframe=teal!35!black,title={Before you move on}]
What does \ja{さめる} mean? (\textquotedblleft{}To wake up; to come round\textquotedblright{}.) Say it once more.
\end{tcolorbox}
\section[\texorpdfstring{\ja{おこす} (okosu)}{おこす (okosu)}]{\ja{おこす} (okosu) — to wake someone; to raise}

\label{lesson:JA-C146-okosu}



\begin{quote}
Before the new one: say the Japanese for soundly, then the Japanese for to wake up.
\end{quote}

\subsection*{You'll want to know: \ja{おこす}}
\textbf{\ja{おこす}} — \emph{okosu} — \textquotedblleft{}to wake someone; to raise\textquotedblright{}.

\begin{quote}
\textbf{\ja{ははは} \ja{まいあさ} \ja{こどもを} \ja{おこします。}} — \emph{haha wa maiasa kodomo o okoshimasu} — \textquotedblleft{}My mother wakes the children every morning.\textquotedblright{}
\end{quote}

Polite \textbf{\ja{おこします}}, \emph{okoshimasu}; a godan verb. Its intransitive partner is \textbf{\ja{おきます}}, \emph{okimasu}.

\begin{cousinweb}
\textbf{\ja{おこす}} (\emph{okosu}) is a native Japanese word (wago), inherited rather than borrowed from Chinese, and built on the same root as \emph{okiru}, \textquotedblleft{}to get up\textquotedblright{}; its kanji was attached later as a kun'yomi (JMdict, KANJIDIC2).
\end{cousinweb}

\begin{grammarlens}[title={What you've built}]
One more word for everyday life at A2.
\end{grammarlens}

\subsection*{Your turn}
\begin{itemize}
\raggedright
  \item \emph{Say it:} \emph{okosu}
  \item \emph{Say it:} \emph{okosu}, once more
  \item \emph{Say it:} \emph{haha wa maiasa kodomo o okoshimasu} — My mother wakes the children every morning.
  \item \emph{Recall:} say the Japanese for soundly, then the Japanese for to wake up, then say \emph{okosu} again
\end{itemize}

\begin{tcolorbox}[breakable,colback=teal!4,colframe=teal!35!black,title={Before you move on}]
What does \ja{おこす} mean? (\textquotedblleft{}To wake someone; to raise\textquotedblright{}.) Say it once more.
\end{tcolorbox}
\section[\texorpdfstring{\ja{すいみん} (suimin)}{すいみん (suimin)}]{\ja{すいみん} (suimin) — sleep}

\label{lesson:JA-C146-suimin}



\begin{quote}
Before the new one: say the Japanese for to wake up, then the Japanese for to wake someone.
\end{quote}

\subsection*{You'll want to know: \ja{すいみん}}
\textbf{\ja{すいみん}} — \emph{suimin} — \textquotedblleft{}sleep\textquotedblright{}.

\begin{quote}
\textbf{\ja{すいみんは} \ja{たいせつです。}} — \emph{suimin wa taisetsu desu} — \textquotedblleft{}Sleep is important.\textquotedblright{}
\end{quote}

\begin{cousinweb}
\textbf{\ja{すいみん}} (\emph{suimin}) is Sino-Japanese (kango): two Chinese characters read with their on'yomi, readings borrowed from Chinese. They are \emph{sui}, \textquotedblleft{}drowsy, sleep\textquotedblright{} (Middle Chinese \emph{dzyweH}) and \emph{min}, \textquotedblleft{}sleep\textquotedblright{} (\emph{men}) (KANJIDIC2; Middle Chinese in Baxter's transcription).
\end{cousinweb}

\begin{grammarlens}[title={What you've built}]
One more word for everyday life at A2.
\end{grammarlens}

\subsection*{Your turn}
\begin{itemize}
\raggedright
  \item \emph{Say it:} \emph{suimin}
  \item \emph{Say it:} \emph{suimin}, once more
  \item \emph{Say it:} \emph{suimin wa taisetsu desu} — Sleep is important.
  \item \emph{Recall:} say the Japanese for to wake up, then the Japanese for to wake someone, then say \emph{suimin} again
\end{itemize}

\begin{tcolorbox}[breakable,colback=teal!4,colframe=teal!35!black,title={Before you move on}]
What does \ja{すいみん} mean? (\textquotedblleft{}Sleep\textquotedblright{}.) Say it once more.
\end{tcolorbox}
